% Image credits for the photographs and AI illustrations of Book 4
% (University Biology, Year 2). Machine-readable license records live in
% images/book4/CREDITS.md; keep the two files in sync.
\chapter*{चित्र-श्रेय}

% Under bidi=basic a Latin run is ONE directional box, so a long author or
% institution name cannot be broken at all; \allowbreak inside it does nothing.
% The Arabic edition reported this page as the book's only overfull boxes.
% \sloppy must be in force when the paragraph is BROKEN, hence the \par before
% \endgroup -- without it the group closes first and the tolerance is restored.
\begingroup\sloppy

इस पुस्तक के रेखाचित्र मौलिक हैं। छायाचित्र अपने-अपने मुक्त अनुज्ञापत्रों के
अंतर्गत प्रयुक्त हैं; उनके रचयिताओं के प्रति आभार:

\begin{itemize}
  \item अध्याय~\ref{ch:b2:microbial-metabolism}: Sergei Winogradsky का
        छायाचित्र, 1890 का दशक --- सार्वजनिक अधिकार-क्षेत्र
        (Wikimedia Commons)।
  \item अध्याय~\ref{ch:b2:genome-diversification}: अपनी प्रयोगशाला में
        Barbara McClintock का छायाचित्र, Smithsonian Institution /
        Science Service, 1947 --- सार्वजनिक अधिकार-क्षेत्र
        (Wikimedia Commons)।
  \item अध्याय~\ref{ch:b2:blood-circulation}: William Harvey का चित्र,
        और उनकी \emph{Exercitatio anatomica de motu cordis} (1628) से
        अग्रबाहु की शिराओं की पट्टिका, Wellcome Collection ---
        CC~BY~4.0 (Wikimedia Commons); एक लाल रक्त कोशिका, एक बिंबाणु
        और एक लसीकाणु का क्रमवीक्षी इलेक्ट्रॉन सूक्ष्मचित्र, Electron
        Microscopy Facility, National Cancer Institute --- सार्वजनिक
        अधिकार-क्षेत्र (Wikimedia Commons)।
  \item अध्याय~\ref{ch:b2:speciation}: John Gould द्वारा बनाई डार्विन
        की फ़िंचों की पट्टिका, \emph{Journal of Researches} (1845) से
        --- सार्वजनिक अधिकार-क्षेत्र (Wikimedia Commons)।
  \item अध्याय~\ref{ch:b2:biogeochemical-cycles}: Mauna Loa वेधशाला
        (Forrest M. Mims III, 2006) और वहाँ लिया जाता एक फ़्लास्क-प्रतिदर्श
        (Commander John Bortniak, 1982), NOAA Photo Library ---
        सार्वजनिक अधिकार-क्षेत्र (Wikimedia Commons)।
\end{itemize}

प्रत्येक छायाचित्र के पूर्ण स्रोत-लिंक और अनुज्ञापत्र के पते पुस्तक के
भंडार में \texttt{images/\allowbreak book4/\allowbreak CREDITS.md} में दिए गए हैं।

\bigskip
छायाचित्रों और मौलिक रेखाचित्रों के साथ-साथ, इस खंड के अध्यायों में फैले
लगभग सत्तर चित्र कृत्रिम बुद्धिमत्ता से बनी चित्रांकन-कृतियाँ हैं, जो
पुस्तक के संपादकों द्वारा लिखे संकेतों से OpenAI की चित्र-जनन प्रणाली से
बनाई गईं और सम्मिलित करने से पहले जैविक शुद्धता के लिए जाँची गईं।
बनाए गए हर चित्र का पूरा संकेत इस किताब के भंडार में, फ़ाइल
\texttt{images/\allowbreak book4/\allowbreak ai/\allowbreak PROMPTS.md} में दर्ज है।
\par\endgroup
\clearpage
